\documentclass[../../main.tex]{subfiles}

\graphicspath{{../../image/}{../../figure/}}

\begin{document}

\section{习题}\label{section:ch05-习题}

\begin{exercise}\label{exercise:ch05-01}
（$\star$）考虑一个涉及$K$个类和目标向量$\boldsymbol{t}$的分类问题，使用1-of-$K$二进制编码方案。证明条件期望$\mathbb{E}[\boldsymbol{t}\mid\boldsymbol{x}]$由后验概率$p(C_k\mid\boldsymbol{x})$给出。
\end{exercise}

\begin{exercise}\label{exercise:ch05-02}
（$\star\star$）在给定一组数据点$\{\boldsymbol{x}_n\}$的情况下，我们可以将凸包定义为所有点$\boldsymbol{x}$的集合
\begin{align}
  \boldsymbol{x}=\sum_{n}\alpha_n\boldsymbol{x}_n
  \label{align:ch05-凸包}
\end{align}
其中$\alpha_n\geqslant 0$且$\sum_n\alpha_n=1$。考虑另一组数据点$\{\boldsymbol{y}_n\}$及相应的凸包。根据定义，如果存在一个向量$\boldsymbol{w}$和一个标量$w_0$，使得对所有$\boldsymbol{x}_n$而言$\boldsymbol{w}^{\mathsf{T}}\boldsymbol{x}_n+w_0>0$，并且对所有$\boldsymbol{y}_n$而言$\boldsymbol{w}^{\mathsf{T}}\boldsymbol{y}_n+w_0<0$，那么这两组数据点就是线性可分的。证明如果它们的凸包相交，则这两组数据点不可能是线性可分的；反之，如果它们是线性可分的，则它们的凸包不相交。
\end{exercise}

\begin{exercise}\label{exercise:ch05-03}
（$\star\star$）考虑最小化平方和误差函数\eqref{align:ch05-平方和分类误差}，并假设训练集中所有的目标向量都满足线性约束条件
\begin{align}
  \boldsymbol{a}^{\mathsf{T}}\boldsymbol{t}_n+b=0
  \label{align:ch05-习题约束}
\end{align}
其中$\boldsymbol{t}_n$相当于\eqref{align:ch05-平方和分类误差}中矩阵$\boldsymbol{T}$的第$n$行。证明作为以上约束条件的结果，由最小二乘解\eqref{align:ch05-分类判别函数}给出的模型预测$\boldsymbol{y}(\boldsymbol{x})$的元素也满足该约束条件，即
\begin{align}
  \boldsymbol{a}^{\mathsf{T}}\boldsymbol{y}(\boldsymbol{x})+b=0
  \label{align:ch05-习题预测约束}
\end{align}
为此，假设其中一个基函数$\phi_0(\boldsymbol{x})=1$，这样相应的参数$w_0$就起到了偏置的作用。
\end{exercise}

\begin{exercise}\label{exercise:ch05-04}
（$\star\star$）扩展\cref{exercise:ch05-03}的结果，证明如果目标向量同时满足多个线性约束条件，那么线性模型的最小二乘预测也将满足相同的约束条件。
\end{exercise}

\begin{exercise}\label{exercise:ch05-05}
（$\star$）利用定义\eqref{align:ch05-F分数}以及\eqref{align:ch05-精确率}和\eqref{align:ch05-召回率}推导$F$分数的结果\eqref{align:ch05-F分数展开}。
\end{exercise}

\begin{exercise}\label{exercise:ch05-06}
（$\star\star$）考虑两个非负数$a$和$b$，并证明如果$a\leqslant b$，则$a\leqslant(ab)^{1/2}$。利用这一结果可以证明，如果选择二分类问题的决策区域是为了最小化误分类概率，则该概率将满足
\begin{align}
  p(\text{错误})\leqslant\int\bigl\{p(\boldsymbol{x},C_1)\,p(\boldsymbol{x},C_2)\bigr\}^{1/2}\mathrm{d}\boldsymbol{x}
  \label{align:ch05-习题错误界}
\end{align}
\end{exercise}

\begin{exercise}\label{exercise:ch05-07}
（$\star$）给定一个含有元素$L_{kj}$的损失矩阵，如果我们为每个$\boldsymbol{x}$选择最小化\eqref{align:ch05-最小风险}的类，则预期风险最小。验证一下，当损失矩阵为$L_{kj}=1-I_{kj}$（其中$I_{kj}$为单位矩阵的元素）时，则可简化为选择后验概率最大的类的准则。如何解释这种形式的损失矩阵？
\end{exercise}

\begin{exercise}\label{exercise:ch05-08}
（$\star$）在有关的一般损失矩阵和一般先验概率时，推导最小化预期损失的准则。
\end{exercise}

\begin{exercise}\label{exercise:ch05-09}
（$\star$）考虑一组$N$个数据点的后验概率的平均值，其形式为
\begin{align}
  \frac{1}{N}\sum_{n=1}^{N}p(C_k\mid\boldsymbol{x}_n)
  \label{align:ch05-习题后验平均}
\end{align}
通过取极限$N\to\infty$，证明这个量接近类的先验概率$p(C_k)$。
\end{exercise}

\begin{exercise}\label{exercise:ch05-10}
（$\star\star$）考虑一个分类问题，其中当来自$C_k$类的输入向量被归入$C_j$类时，产生的损失由损失矩阵$L_{kj}$给出，而选择拒绝选项所产生的损失为$\lambda$。推导预期损失最小的决策准则。验证一下，当损失矩阵由$L_{kj}=1-I_{kj}$给出时，所要推导的决策准则可简化为\cref{subsection:ch05-拒绝选项}讨论的拒绝准则。$\lambda$和拒绝阈值$\theta$之间有何关系？
\end{exercise}

\begin{exercise}\label{exercise:ch05-11}
（$\star$）证明 logistic sigmoid 函数\,[即\eqref{align:ch05-sigmoid定义}]\,满足性质$\sigma(-a)=1-\sigma(a)$，其逆函数由$\sigma^{-1}(y)=\ln\{y/(1-y)\}$给出。
\end{exercise}

\begin{exercise}\label{exercise:ch05-12}
（$\star$）利用\eqref{align:ch05-二类后验sigmoid}和\eqref{align:ch05-预激活对数比}得出具有 Gauss 密度的二类生成式模型的类的后验概率结果\eqref{align:ch05-Gauss后验}，并对参数$\boldsymbol{w}$和$w_0$的结果\eqref{align:ch05-Gauss权重}和\eqref{align:ch05-Gauss偏置}进行验证。
\end{exercise}

\begin{exercise}\label{exercise:ch05-13}
（$\star$）考虑一个由类的先验概率$p(C_k)=\pi_k$和类-条件概率密度$p(\boldsymbol{\phi}\mid C_k)$定义的涉及$K$个类的分类模型，其中$\boldsymbol{\phi}$是输入特征。假设我们得到一个训练数据集$\{\boldsymbol{\phi}_n,\boldsymbol{t}_n\}$，其中$n=1,\cdots,N$，$\boldsymbol{t}_n$是长度为$K$且采用1-of-$K$编码方案的二进制目标向量。如果数据点$n$属于$C_k$类，则$t_{nj}=I_{jk}$。假设数据点是从该模型中独立提取的，证明先验概率的最大似然解为
\begin{align}
  \pi_k=\frac{N_k}{N}
  \label{align:ch05-习题先验ML}
\end{align}
其中$N_k$是归入$C_k$类的数据点数量。
\end{exercise}

\begin{exercise}\label{exercise:ch05-14}
（$\star\star$）考虑\cref{exercise:ch05-13}中的分类模型，假设类-条件概率密度由具有共享协方差矩阵的 Gauss 分布给出，因此
\begin{align}
  p(\boldsymbol{\phi}\mid C_k)=\mathcal{N}(\boldsymbol{\phi}\mid\boldsymbol{\mu}_k,\boldsymbol{\Sigma})
  \label{align:ch05-习题Gauss}
\end{align}
证明$C_k$类的 Gauss 分布均值的最大似然解由下式给出：
\begin{align}
  \boldsymbol{\mu}_k=\frac{1}{N_k}\sum_{n=1}^{N}t_{nk}\boldsymbol{\phi}_n
  \label{align:ch05-习题均值ML}
\end{align}
它表示归入$C_k$类的那些特征向量的均值。类似地，证明共享协方差矩阵的最大似然解由下式给出：
\begin{align}
  \boldsymbol{\Sigma}=\sum_{k=1}^{K}\frac{N_k}{N}\boldsymbol{S}_k
  \label{align:ch05-习题协方差ML}
\end{align}
其中
\begin{align}
  \boldsymbol{S}_k=\frac{1}{N_k}\sum_{n=1}^{N}t_{nk}(\boldsymbol{\phi}_n-\boldsymbol{\mu}_k)(\boldsymbol{\phi}_n-\boldsymbol{\mu}_k)^{\mathsf{T}}
  \label{align:ch05-习题类协方差}
\end{align}
因此，$\boldsymbol{\Sigma}$由与每个类相关的数据协方差的加权平均值给出，其中加权系数由类的先验概率给出。
\end{exercise}

\begin{exercise}\label{exercise:ch05-15}
（$\star\star$）推导\cref{subsection:ch05-离散特征}所讲的具有离散二元特征的概率朴素 Bayes 分类器参数$\{\mu_{ki}\}$的最大似然解。
\end{exercise}

\begin{exercise}\label{exercise:ch05-16}
（$\star\star$）考虑一个涉及$K$个类的分类问题，其中输入特征$\boldsymbol{\phi}$有$M$个组件，每个组件有$L$个离散状态。将组件的值用1-of-$L$二进制编码方案表示。进一步假设，以$C_k$类为条件，$\boldsymbol{\phi}$的$M$个组件是独立的，因此关于特征向量组件的类-条件概率密度可以分解。

证明由\eqref{align:ch05-预激活定义}给出的量$a_k$（出现在描述类的后验概率的 softmax 函数的自变量中）是$\boldsymbol{\phi}$的组件的线性函数。请注意，这是一个朴素 Bayes 模型（参见\cref{subsection:ch11-朴素Bayes}）的示例。
\end{exercise}

\begin{exercise}\label{exercise:ch05-17}
（$\star\star$）推导\cref{exercise:ch05-16}中描述的概率朴素 Bayes 分类器参数的最大似然解。
\end{exercise}

\begin{exercise}\label{exercise:ch05-18}
（$\star$）验证由\eqref{align:ch05-sigmoid定义}定义的 logistic sigmoid 函数导数的关系式\eqref{align:ch05-sigmoid导数}。
\end{exercise}

\begin{exercise}\label{exercise:ch05-19}
（$\star$）利用 logistic sigmoid 函数的导数结果\eqref{align:ch05-sigmoid导数}，证明 logistic 回归模型的误差函数\eqref{align:ch05-交叉熵}由\eqref{align:ch05-交叉熵梯度式}给出。
\end{exercise}

\begin{exercise}\label{exercise:ch05-20}
（$\star$）证明对于线性可分数据集，logistic 回归模型的最大似然解可以通过计算一个向量$\boldsymbol{w}$来获得，该向量的决策边界$\boldsymbol{w}^{\mathsf{T}}\boldsymbol{\phi}(\boldsymbol{x})=0$可将类分开，然后将$\boldsymbol{w}$的大小取到无穷大。
\end{exercise}

\begin{exercise}\label{exercise:ch05-21}
（$\star$）证明 softmax 激活函数\,[即\eqref{align:ch05-多类后验}]\,的导数\,[其中$a_k$由\eqref{align:ch05-多类预激活}定义]\,由\eqref{align:ch05-softmax导数式}给出。
\end{exercise}

\begin{exercise}\label{exercise:ch05-22}
（$\star$）利用 softmax 激活函数的导数结果\,[即\eqref{align:ch05-softmax导数式}]，证明交叉熵误差\,[即\eqref{align:ch05-多类交叉熵}]\,的梯度由\eqref{align:ch05-多类梯度式}给出。
\end{exercise}

\begin{exercise}\label{exercise:ch05-23}
（$\star$）证明 probit 函数\,[即\eqref{align:ch05-probit函数}]\,和 erf 函数\,[即\eqref{align:ch05-erf函数}]\,的关系式\eqref{align:ch05-probit-erf}。
\end{exercise}

\begin{exercise}\label{exercise:ch05-24}
（$\star\star$）假设我们希望用缩放 probit 函数$\Phi(\lambda a)$来近似由\eqref{align:ch05-sigmoid定义}定义的 logistic sigmoid 函数$\sigma(a)$，其中$\Phi(a)$由\eqref{align:ch05-probit函数}定义。证明如果选择的$\lambda$能使这两个函数的导数在$a=0$处相等，则$\lambda^{2}=\pi/8$。
\end{exercise}

\end{document}
